\section{Introduction}

% Q1: What is the problem and why do I care?
Solving imperfect-information games is a major challenge in multiagent systems, with
applications from poker to security~\cite{tambe2011security}, where agents
have to make strategic decisions without being able to observe the complete state of the
environment or the private information available to their opponents. A representative example
of this is poker, where players have to choose actions based on their private cards, public
information, the betting history, and their beliefs about other players' hidden
cards~\cite{bowling2015heads, moravcik2017deepstack, brown2018superhuman, brown2019superhuman}. A main
goal in these settings is to calculate strategies that are difficult for the opponent to
exploit. A widely used approach for this is Counterfactual Regret Minimization
(CFR)~\cite{zinkevich2007regret}, since it iteratively improves strategies by minimizing the
regret over possible actions at every information set. However, convergence may require many
iterations, making computational efficiency an important concern when solving larger
games~\cite{tammelin2015solving}.

% Q2: Why is it important/difficult?
It is difficult to solve imperfect-information games, since an agent has to account for
uncertainty regarding the underlying game state while also reasoning about the strategic
responses of other agents. Unlike in perfect-information games, an agent cannot select actions
based on a fully observed state. It must, instead, learn a strategy that will perform well
across the multiple possible states that are consistent with its observations. CFR addresses
this challenge by decomposing regret minimization across information sets; however, the amount
and pattern of regret may vary substantially across these sets. Some information sets might
exhibit rapidly changing regret, while others might have relatively stable regret. Applying the
same update procedure throughout the game might therefore fail to exploit the differences in
local learning dynamics.

% Q3: What has been done already in this problem area?
Several methods have already improved the efficiency of regret minimization. CFR introduced an
iterative framework for minimizing counterfactual regret and established convergence guarantees
for the average strategy in two-player zero-sum games~\cite{zinkevich2007regret}. \cfrp{} then
improved practical convergence by using regret matching$^+$, alternating player updates, and
linear averaging~\cite{tammelin2015solving}. Discounted CFR introduced discounting schemes that
reduce the influence of earlier iterations and can accelerate convergence~\cite{brown2019solving}.
Predictive \cfrp{} further explores the use of predicted regret updates to improve game-solving
speed~\cite{farina2021faster}, and PDCFR$^+$ combines prediction with
discounting~\cite{xu2024minimizing}. More recent methods also adapt these choices automatically:
DDCFR learns its discounting with reinforcement learning~\cite{xu2024dynamic}, hyperparameter
schedules change the discounting over time~\cite{zhang2026faster}, VAD-CFR adjusts discounting
using a single, game-wide measure of regret volatility~\cite{li2026discovering}, and APCFR$^+$
adapts the size of its predicted regret updates separately at each information
set~\cite{meng2025faster}. Together, these methods demonstrate that how regret is modified,
accumulated, weighted, and updated can drastically affect convergence.

% Q4: What particular problem remains unsolved?
Despite all of these advances, the question of whether the discounting of past regret should
adapt to the recent dynamics of individual information sets remains worth investigating. It is
especially worth noting that a fixed discounting procedure might not respond equally well to
information sets whose regrets are changing rapidly and to those whose regrets have stabilized.
Even though existing approaches have modified regret updates, used predictive information, or
adapted their discounting for the game as a whole, this work investigates a complementary
approach: adapting the regret update locally according to recent regret dynamics, and whether
this improves convergence. An open empirical question is under which conditions
this would be most beneficial.

% Q5: How did you solve it?
In this work, we present two methods. The first, Locally Adaptive Discounted \cfrp{} (\ours{}),
is a modification of \cfrp{} that uses the recent regret history at individual information
sets to determine how much past regret should be kept. Instead of using an identical discount
at each information set, it adjusts the discount according to the observed local regret
dynamics. The second, \mix{}, goes further: each information set runs \cfrp{}, Discounted CFR
and Predictive \cfrp{} side by side, and a small regret-matching learner at that information
set learns which of them to follow. Both algorithms are implemented using
OpenSpiel~\cite{lanctot2019openspiel} and evaluated on Kuhn Poker~\cite{kuhn1950simplified} and
Leduc Poker~\cite{southey2005bayes}, two imperfect-information poker games that provide
increasingly complex test environments, and on three further games, two sizes of Liar's dice
and Goofspiel, to test whether the results transfer beyond poker.

% Q6: What is cool about your approach?
Our main motivation behind \ours{} is that different parts of a game might exhibit different
learning dynamics. By incorporating the local regret history into the update procedure, the
method aims to discount past regret more strongly at information sets where the regret is
changing greatly, so that they respond quickly, while keeping the regret history where it has
stabilized. \mix{} applies the same idea to the choice of update rule itself: since no existing variant
is fastest on every game, each information set can follow whichever rule works best there,
and regret matching over the rules keeps the method no-regret without any parameter to tune.
In contrast to a uniform modification applied throughout the game, our approach explicitly
examines the variation in regret behavior across information sets.

% Q7: What were your key results?
To evaluate our methods, we compare them against \cfrp{}, Discounted CFR, and Predictive
\cfrp{}. As a first step, we reproduce the CFR versus \cfrp{} comparison on Kuhn and Leduc
Poker using OpenSpiel, whose CFR implementation uses alternating updates. \cfrp{} reaches
roughly $12\times$ lower exploitability than CFR on Kuhn Poker after 2000 iterations and
$46\times$ lower on Leduc Poker after 1000 iterations. Our main measure is exploitability,
which is how much a best-responding opponent could win against the agents' average strategy,
as a function of the number of iterations. We also compare each method with an otherwise
identical version that adapts at the level of the whole game, to test whether adapting locally
is what makes the difference. We find that no existing variant is fastest on every game:
Predictive \cfrp{} is best on Kuhn Poker but worst on Leduc Poker. \ours{} improves on \cfrp{}
with the same averaging on four of five games, and its local version beats its game-wide
version on all five, but its gains are modest. \mix{}, without any tuning, is the best method
on Kuhn and Leduc Poker and within $1.13\times$ of the best on the other three games, and on
two games it outperforms every rule it combines.

% Q8: What are the contributions of this paper?
The main contributions of this work are:
\begin{enumerate}
  \item We develop an adaptive regret-update approach, \ours{}, that incorporates local
        regret dynamics into \cfrp{} by adapting how much past regret each information set
        keeps.
  \item We develop \mix{}, which selects among existing CFR update rules separately at each
        information set, and show that it remains no-regret.
  \item We establish an experimental comparison of our methods against established
        regret-minimization baselines on Kuhn Poker, Leduc Poker, Liar's dice and Goofspiel.
  \item We investigate the conditions under which local adaptation improves
        convergence, including differences between game environments and the effects of
        individual components through ablation studies.
\end{enumerate}
The results provide evidence regarding the potential benefits and limitations of adapting
regret updates at the information-set level.
